\documentclass[aspectratio=169, 10pt]{beamer}

% ==========================================
% TEMA Y CONFIGURACIÓN VISUAL
% ==========================================
\usetheme{Boadilla}
\usecolortheme{default}
\setbeamertemplate{navigation symbols}{} % Quita los iconos de navegación inferiores
\setbeamertemplate{itemize items}[circle]
\setbeamertemplate{enumerate items}[circle]
\setbeamertemplate{blocks}[rounded][shadow=false]

% ==========================================
% PAQUETES
% ==========================================
\usepackage[utf8]{inputenc}
\usepackage[T1]{fontenc}
\usepackage[spanish]{babel}
\usepackage{graphicx}
\usepackage{booktabs}
\usepackage{multicol}
\usepackage{tikz}
\usepackage{pgfplots}
\usepackage{fontawesome5}
\usepackage{hyperref}
\usetikzlibrary{positioning, shapes.geometric, arrows.meta, calc}
\pgfplotsset{compat=1.18}

% ==========================================
% COLORES CORPORATIVOS PROESA
% ==========================================
\definecolor{proesadark}{RGB}{5,28,44}      % Azul oscuro casi negro
\definecolor{proesablue}{RGB}{0,86,179}     % Azul vibrante
\definecolor{proesalight}{RGB}{0,164,228}   % Celeste
\definecolor{proesagray}{RGB}{243,244,246}  % Gris claro para fondos
\definecolor{proesagreen}{RGB}{22,163,74}   % Verde estado/éxito

% Asignación de colores a elementos de Beamer
\setbeamercolor{structure}{fg=proesablue}
\setbeamercolor{title}{fg=proesadark}
\setbeamercolor{frametitle}{bg=proesadark, fg=white}
\setbeamercolor{normal text}{fg=darkgray}
\setbeamercolor{block title}{bg=proesablue, fg=white}
\setbeamercolor{block body}{bg=proesagray, fg=darkgray}
\setbeamercolor{alertblock title}{bg=proesalight, fg=white}
\setbeamercolor{alertblock body}{bg=proesagray, fg=darkgray}

% ==========================================
% METADATOS
% ==========================================
\title{\textbf{Presentación Final del Proyecto}}
\subtitle{Sistema de Inteligencia de Mercado: Webscraping, Estandarización MDM y Auditoría INVIMA de Alcohol y Tabaco}
\author{PROESA – Centro de Estudios en Protección Social y Economía de la Salud}
\institute{Universidad Icesi | Banco Mundial}
\date{Agosto de 2026}

% ==========================================
% DOCUMENTO
% ==========================================
\begin{document}

% --- DIAPOSITIVA 1: PORTADA ---
\begin{frame}[plain]
    \begin{tikzpicture}[remember picture,overlay]
        % Barra decorativa superior
        \fill[proesablue] (current page.north west) rectangle ([yshift=-0.18cm]current page.north east);
        % Fondo inferior sutil
        \fill[proesagray] (current page.south west) rectangle ([yshift=1.8cm]current page.south east);
    \end{tikzpicture}
    
    \vspace{0.8cm}
    {\color{proesadark}\Huge \textbf{Informe Final de Resultados}} \\
    \vspace{0.3cm}
    {\color{proesablue}\Large Webscraping para Alcohol y Tabaco en Colombia}
    
    \vspace{1.2cm}
    {\small \textbf{Producto Final} \\ Presentación Ejecutiva}
    
    \vfill
    {\small \textbf{PROESA – Universidad Icesi} \\ \href{http://www.icesi.edu.co/proesa}{www.icesi.edu.co/proesa} \\ Agosto de 2026}
\end{frame}

% --- DIAPOSITIVA 2: RESUMEN EJECUTIVO DE ENTREGABLES ---
\begin{frame}{Resumen Ejecutivo de Entregables}
    \begin{itemize}
        \setlength{\itemsep}{10pt}
        \item \textbf{Arquitectura Automatizada End-to-End:} Extracción diaria, limpia y resiliente sobre 7 grandes cadenas de supermercados y comercios digitales en Colombia (Éxito, Carulla, Jumbo, D1, Cañaveral, Olímpica, Makro).
        \item \textbf{Maestro Único MDM con IA (V4 Flash con Pensamiento):} Estandarización inteligente de productos mediante modelos de lenguaje con razonamiento CoT y RAG sobre 6,631 históricos locales.
        \item \textbf{Auditoría y Vinculación INVIMA:} Asignación automatizada de Registros Sanitarios oficiales desde el catálogo INVIMA con política estricta de cero falsos positivos.
        \item \textbf{Suite Gráfica \& Alertas Automáticas:} Interfaz gráfica integral en Python/CustomTkinter (\texttt{suite\_app.py}) y sistema de notificaciones ejecutivas por correo con Resend.
    \end{itemize}
\end{frame}

% --- DIAPOSITIVA 3: ARQUITECTURA TÉCNICA DEL SISTEMA ---
\begin{frame}{Arquitectura del Pipeline de Inteligencia de Mercado}
    \begin{center}
        \begin{tikzpicture}[scale=0.88, every node/.style={transform shape}]
            % Fila Superior
            \node (step1) [draw=proesablue, fill=proesagray, rounded corners, thick, text width=3.6cm, align=center, inner sep=6pt] at (0, 2)
                {\textbf{\faDatabase\ 1. Scrapers (7 Comercios)}\\[2pt]\scriptsize Extracción masiva diaria Éxito, Carulla, Jumbo, D1, Cañaveral, Olímpica, Makro};
            
            \node (step2) [draw=proesadark, fill=white, rounded corners, thick, text width=3.6cm, align=center, inner sep=6pt] at (4.8, 2)
                {\textbf{\faDatabase\ 2. ETL Inicial SQLite}\\[2pt]\scriptsize Mapeo automático de conocidos en $<0.5$s (90-95\% catálogo)};
            
            \node (step3) [draw=proesablue, fill=proesablue!10, rounded corners, thick, text width=3.6cm, align=center, inner sep=6pt] at (9.6, 2)
                {\textbf{\faRobot\ 3. IA MDM Matcher}\\[2pt]\scriptsize Modelo V4 Flash con Pensamiento (Razonamiento CoT + RAG)};

            % Fila Inferior
            \node (step6) [draw=proesalight, fill=white, rounded corners, thick, text width=3.6cm, align=center, inner sep=6pt] at (0, -0.6)
                {\textbf{\faEnvelope\ 6. Alertas Resend}\\[2pt]\scriptsize Reporte HTML diario consolidado en bandeja de entrada};

            \node (step5) [draw=proesadark, fill=proesagray, rounded corners, thick, text width=3.6cm, align=center, inner sep=6pt] at (4.8, -0.6)
                {\textbf{\faSync\ 5. ETL Final \& Export}\\[2pt]\scriptsize Vista normalizada \& respaldo \texttt{mdm\_export.json}};

            \node (step4) [draw=proesagreen, fill=proesagreen!10, rounded corners, thick, text width=3.6cm, align=center, inner sep=6pt] at (9.6, -0.6)
                {\textbf{\faFile\ 4. Asignación INVIMA}\\[2pt]\scriptsize Auditoría de certificados oficiales sanitarios};

            % Conexiones
            \draw[->, >=stealth, ultra thick, proesablue] (step1) -- (step2);
            \draw[->, >=stealth, ultra thick, proesablue] (step2) -- (step3);
            \draw[->, >=stealth, ultra thick, proesablue] (step3) -- (step4);
            \draw[->, >=stealth, ultra thick, proesagreen] (step4) -- (step5);
            \draw[->, >=stealth, ultra thick, proesadark] (step5) -- (step6);
        \end{tikzpicture}
    \end{center}
\end{frame}

% --- DIAPOSITIVA 4: INDICADORES CLAVE DE ESTANDARIZACIÓN ---
\begin{frame}{Resultados de Estandarización MDM e INVIMA}
    \vspace{0.2cm}
    \begin{columns}[t]
        \begin{column}{0.25\textwidth}
            \centering
            {\color{proesablue}\Huge \faBoxes} \\ \vspace{0.2cm}
            {\color{proesadark}\Huge \textbf{2,810}} \\ \vspace{0.2cm}
            \small Productos Únicos Maestros (MDM)
        \end{column}
        \begin{column}{0.25\textwidth}
            \centering
            {\color{proesalight}\Huge \faLink} \\ \vspace{0.2cm}
            {\color{proesadark}\Huge \textbf{7,712}} \\ \vspace{0.2cm}
            \small Mapeos Tienda $\rightarrow$ Maestro Estandarizados
        \end{column}
        \begin{column}{0.25\textwidth}
            \centering
            {\color{proesagreen}\Huge \faCheck} \\ \vspace{0.2cm}
            {\color{proesadark}\Huge \textbf{2,128}} \\ \vspace{0.2cm}
            \small Registros Sanitarios INVIMA Asignados (75.7\%)
        \end{column}
        \begin{column}{0.25\textwidth}
            \centering
            {\color{proesablue}\Huge \faFile} \\ \vspace{0.2cm}
            {\color{proesadark}\Huge \textbf{10,972}} \\ \vspace{0.2cm}
            \small Certificados INVIMA Ingestados en BD
        \end{column}
    \end{columns}
    \vspace{0.6cm}
    \begin{block}{Garantía de Consistencia Estandarizada}
        Cada producto extraído diariamente se homologa contra el catálogo \textbf{Maestro MDM}, permitiendo comparar el mismo ítem (ej. \textit{Aguardiente Antioqueño Sin Azúcar 750 ml}) a través de múltiples comercios y variaciones de precio sin duplicidad de referencias.
    \end{block}
\end{frame}

% --- DIAPOSITIVA 5: COBERTURA COMERCIAL ---
\begin{frame}{Cobertura de Cadenas Comerciales Evaluadas}
    \begin{columns}[c]
        \begin{column}{0.48\textwidth}
            \begin{block}{\faStore\ Cadenas Monitoreadas Diariamente}
                \begin{itemize}
                    \setlength{\itemsep}{6pt}
                    \item \textbf{Éxito}
                    \item \textbf{Carulla}
                    \item \textbf{Jumbo}
                    \item \textbf{D1}
                    \item \textbf{Cañaveral}
                    \item \textbf{Olímpica}
                    \item \textbf{Makro}
                \end{itemize}
            \end{block}
        \end{column}
        
        \begin{column}{0.48\textwidth}
            \begin{alertblock}{\faLock\ Resiliencia Antibot (Scrapling)}
                Todos los scrapers operan bajo el framework \textbf{Scrapling} en modo \textbf{HTTP ONLY}, simulando firmas TLS/Headers para evadir antibots (Cloudflare) con consumo mínimo de memoria RAM en el servidor Windows Server 2016.
            \end{alertblock}
        \end{column}
    \end{columns}
\end{frame}

% --- DIAPOSITIVA: DISTRIBUCIÓN POR COMERCIO (ALCOHOL vs TABACO) ---
\begin{frame}{Distribución de Productos por Comercio (Alcohol vs. Tabaco)}
    \begin{columns}[T]
        % Columna Izquierda: Solo Alcohol
        \begin{column}{0.48\textwidth}
            \begin{block}{\centering \faWineGlass\ Solo Alcohol por Comercio}
                \vspace{0.1cm}
                \begin{center}
                    \begin{tikzpicture}[scale=0.85]
                        % Éxito (1,800)
                        \fill[proesablue, rounded corners=2pt] (0, 3.6) rectangle (4.2, 3.95);
                        \node[left, font=\tiny\bfseries] at (-0.1, 3.775) {Éxito};
                        \node[right, font=\tiny\bfseries, text=proesablue] at (4.25, 3.775) {1,800};

                        % Carulla (1,706)
                        \fill[proesablue!90, rounded corners=2pt] (0, 3.0) rectangle (3.98, 3.35);
                        \node[left, font=\tiny\bfseries] at (-0.1, 3.175) {Carulla};
                        \node[right, font=\tiny\bfseries, text=proesablue] at (4.03, 3.175) {1,706};

                        % Olímpica (1,384)
                        \fill[proesablue!80, rounded corners=2pt] (0, 2.4) rectangle (3.23, 2.75);
                        \node[left, font=\tiny\bfseries] at (-0.1, 2.575) {Olímpica};
                        \node[right, font=\tiny\bfseries, text=proesablue] at (3.28, 2.575) {1,384};

                        % Jumbo (1,331)
                        \fill[proesablue!70, rounded corners=2pt] (0, 1.8) rectangle (3.10, 2.15);
                        \node[left, font=\tiny\bfseries] at (-0.1, 1.975) {Jumbo};
                        \node[right, font=\tiny\bfseries, text=proesablue] at (3.15, 1.975) {1,331};

                        % Cañaveral (501)
                        \fill[proesalight, rounded corners=2pt] (0, 1.2) rectangle (1.17, 1.55);
                        \node[left, font=\tiny\bfseries] at (-0.1, 1.375) {Cañaveral};
                        \node[right, font=\tiny\bfseries, text=proesalight!80!black] at (1.22, 1.375) {501};

                        % Makro (392)
                        \fill[proesalight!80, rounded corners=2pt] (0, 0.6) rectangle (0.91, 0.95);
                        \node[left, font=\tiny\bfseries] at (-0.1, 0.775) {Makro};
                        \node[right, font=\tiny\bfseries, text=proesalight!80!black] at (0.96, 0.775) {392};

                        % D1 (53)
                        \fill[proesadark!60, rounded corners=2pt] (0, 0.0) rectangle (0.2, 0.35);
                        \node[left, font=\tiny\bfseries] at (-0.1, 0.175) {D1};
                        \node[right, font=\tiny\bfseries, text=proesadark] at (0.25, 0.175) {53};

                        % Eje vertical
                        \draw[gray!50, thick] (0,-0.1) -- (0,4.1);
                    \end{tikzpicture}
                \end{center}
            \end{block}
        \end{column}

        % Columna Derecha: Solo Tabaco
        \begin{column}{0.48\textwidth}
            \begin{block}{\centering \faSmoking\ Solo Tabaco por Comercio}
                \vspace{0.1cm}
                \begin{center}
                    \begin{tikzpicture}[scale=0.85]
                        % Éxito (111)
                        \fill[proesadark, rounded corners=2pt] (0, 3.6) rectangle (4.2, 3.95);
                        \node[left, font=\tiny\bfseries] at (-0.1, 3.775) {Éxito};
                        \node[right, font=\tiny\bfseries, text=proesadark] at (4.25, 3.775) {111};

                        % Carulla (81)
                        \fill[proesadark!85, rounded corners=2pt] (0, 3.0) rectangle (3.06, 3.35);
                        \node[left, font=\tiny\bfseries] at (-0.1, 3.175) {Carulla};
                        \node[right, font=\tiny\bfseries, text=proesadark] at (3.11, 3.175) {81};

                        % Olímpica (44)
                        \fill[proesadark!70, rounded corners=2pt] (0, 2.4) rectangle (1.66, 2.75);
                        \node[left, font=\tiny\bfseries] at (-0.1, 2.575) {Olímpica};
                        \node[right, font=\tiny\bfseries, text=proesadark] at (1.71, 2.575) {44};

                        % Jumbo (44)
                        \fill[proesadark!70, rounded corners=2pt] (0, 1.8) rectangle (1.66, 2.15);
                        \node[left, font=\tiny\bfseries] at (-0.1, 1.975) {Jumbo};
                        \node[right, font=\tiny\bfseries, text=proesadark] at (1.71, 1.975) {44};

                        % Cañaveral (15)
                        \fill[proesadark!55, rounded corners=2pt] (0, 1.2) rectangle (0.57, 1.55);
                        \node[left, font=\tiny\bfseries] at (-0.1, 1.375) {Cañaveral};
                        \node[right, font=\tiny\bfseries, text=proesadark] at (0.62, 1.375) {15};

                        % D1 (3)
                        \fill[proesadark!40, rounded corners=2pt] (0, 0.6) rectangle (0.2, 0.95);
                        \node[left, font=\tiny\bfseries] at (-0.1, 0.775) {D1};
                        \node[right, font=\tiny\bfseries, text=proesadark] at (0.25, 0.775) {3};

                        % Makro (0)
                        \node[left, font=\tiny\bfseries] at (-0.1, 0.175) {Makro};
                        \node[right, font=\tiny\bfseries, text=gray] at (0.05, 0.175) {0};

                        % Eje vertical
                        \draw[gray!50, thick] (0,-0.1) -- (0,4.1);
                    \end{tikzpicture}
                \end{center}
            \end{block}
        \end{column}
    \end{columns}
\end{frame}



% --- DIAPOSITIVA 6: SEGMENTACIÓN DEL CATÁLOGO MAESTRO ---
\begin{frame}{Segmentación del Catálogo Maestro Único}
    \begin{columns}[c]
        \begin{column}{0.45\textwidth}
            \begin{center}
                \begin{tikzpicture}
                    % Gráfica tipo Dona TikZ
                    \fill[proesalight] (0,0) circle (2.3cm);
                    \fill[proesablue] (0,0) -- (90:2.3cm) arc (90:430.56:2.3cm) -- cycle; % 94.6%
                    \fill[white] (0,0) circle (1.4cm);
                    
                    \node[proesadark] at (0,0) {\small \textbf{2,810 Maestros}};
                    \node[white] at (260:1.8cm) {\small \textbf{94.6\%}};
                    \draw[proesadark, thick] (70:2.0cm) -- (70:2.8cm) node[right, text=proesadark] {\small \textbf{5.4\%}};
                \end{tikzpicture}
            \end{center}
        \end{column}

        \begin{column}{0.53\textwidth}
            \begin{block}{\faWineGlass\ Alcohol (2,659 Productos Maestros - 94.6\%)}
                \small Vinos (1,378), Cervezas (363), Whisky (159), Coctelería/Aperitivos (155), Ron (121), Cremas (109), Aguardiente (108), Tequila (89), Ginebra (45).
            \end{block}
            \vspace{0.2cm}
            \begin{block}{\faSmoking\ Tabaco (151 Productos Maestros - 5.4\%)}
                \small Cigarrillos tradicionales, vapeadores, esencias, puros y bolsas de nicotina.
            \end{block}
        \end{column}
    \end{columns}
\end{frame}

% --- DIAPOSITIVA 7: MOTOR DE IA Y RAG HISTÓRICO ---
\begin{frame}{Inteligencia Artificial y Aprendizaje RAG}
    \begin{columns}[T]
        \begin{column}{0.48\textwidth}
            \begin{block}{\faBrain\ Modelo V4 Flash con Pensamiento}
                \begin{itemize}
                    \item \textbf{Razonamiento CoT:} Cadena de pensamiento explícita (\texttt{reasoning\_content}) que analiza volumen, coherencia de precio y variaciones semánticas antes de clasificar.
                    \item \textbf{Alta Precisión:} Evalúa candidatos en tiempo récord resolviendo nombres complejos o ambiguos.
                \end{itemize}
            \end{block}
            \vspace{0.2cm}
            \begin{block}{\faMicrochip\ Memoria de DB Instantánea}
                Índice en memoria de \textbf{4,436 nombres validados} que mapea en 0 ms sin costo de API productos previamente aprobados.
            \end{block}
        \end{column}

        \begin{column}{0.48\textwidth}
            \begin{alertblock}{\faCheck\ Control y Calidad de Datos}
                \begin{itemize}
                    \item \textbf{Falsos Positivos:} Detección y filtrado asistido de accesorios (hieleras, vasos, pasabocas) sin alterar la serie histórica.
                    \item \textbf{Atributos Canónicos:} Estandarización unificada de volumen, marca y graduación alcohólica.
                \end{itemize}
            \end{alertblock}
        \end{column}
    \end{columns}
\end{frame}

% --- DIAPOSITIVA 8: AUDITORÍA SANITARIA INVIMA ---
\begin{frame}{Auditoría de Registros Sanitarios INVIMA}
    \begin{columns}[c]
        \begin{column}{0.5\textwidth}
            \begin{block}{\faFile\ Cobertura del 75.7\% Alcanzada}
                Se logró vincular exitosamente \textbf{2,128 productos maestros de alcohol} con su Registro Sanitario oficial en la tabla \texttt{invima\_certificados}.
            \end{block}
            \vspace{0.3cm}
            \begin{block}{\faCheck\ Criterio Estricto de Coincidencia Sanitaria}
                Si un producto no encuentra una coincidencia exacta de marca y variedad con confianza $\ge 90\%$, la casilla se mantiene en \texttt{NULL}, \textbf{evitando cualquier falso positivo}.
            \end{block}
        \end{column}

        \begin{column}{0.48\textwidth}
            \begin{center}
                \begin{tikzpicture}[scale=0.92]
                    % Ejes
                    \draw[->, >=stealth, gray!70] (0,0) -- (5.2,0);
                    \draw[->, >=stealth, gray!70] (0,0) -- (0,4.0) node[above, text=proesadark, font=\tiny\bfseries] {Cantidad de Maestros};
                    
                    % Rejilla de guía
                    \draw[dashed, gray!30] (0, 3.2) -- (4.8, 3.2) node[right, gray, font=\tiny] {2,000};
                    \draw[dashed, gray!30] (0, 1.6) -- (4.8, 1.6) node[right, gray, font=\tiny] {1,000};

                    % Barra 1: Con INVIMA (2,128)
                    \fill[proesagreen, rounded corners=2pt] (0.5,0) rectangle (1.5, 3.4);
                    \node[above, font=\bfseries\footnotesize, text=proesagreen] at (1.0, 3.4) {2,128};
                    \node[below, font=\scriptsize, align=center] at (1.0, -0.1) {\textbf{Con INVIMA}\\[-2pt]\tiny (75.7\%)};

                    % Barra 2: Pendientes (531)
                    \fill[proesablue, rounded corners=2pt] (2.1,0) rectangle (3.1, 0.85);
                    \node[above, font=\bfseries\footnotesize, text=proesablue] at (2.6, 0.85) {531};
                    \node[below, font=\scriptsize, align=center] at (2.6, -0.1) {\textbf{Pendientes}\\[-2pt]\tiny (18.9\%)};

                    % Barra 3: Tabaco (151)
                    \fill[gray!70, rounded corners=2pt] (3.7,0) rectangle (4.7, 0.24);
                    \node[above, font=\bfseries\footnotesize, text=gray!80!black] at (4.2, 0.24) {151};
                    \node[below, font=\scriptsize, align=center] at (4.2, -0.1) {\textbf{Tabaco}\\[-2pt]\tiny (5.4\%)};
                \end{tikzpicture}
            \end{center}

        \end{column}
    \end{columns}
\end{frame}

% --- DIAPOSITIVA 9: PRÓXIMOS PASOS Y RECOMENDACIONES ---
\begin{frame}{Próximos Pasos y Recomendaciones para Continuidad}
    \begin{enumerate}
        \setlength{\itemsep}{16pt}
        \item \textbf{\color{proesablue}Actualización Periódica de la Base de Datos INVIMA:} 
        El trámite de Registros Sanitarios ante el INVIMA es dinámico (nuevas expediciones, renovaciones y cancelaciones). Para elevar la cobertura del \textbf{75.7\% actual hacia el 100\%}, se recomienda establecer un flujo de actualización periódica consumiendo la base de datos oficial actualizada expedida por la autoridad sanitaria.
        
        \item \textbf{\color{proesadark}Análisis Longitudinal y Elasticidad de Precios:} 
        Aprovechar la serie temporal acumulada en \texttt{productos\_historico} para estudiar elasticidad-precio, dinámicas promocionales y el impacto del marco tributario de impuestos saludables en Colombia.
    \end{enumerate}
\end{frame}

% --- DIAPOSITIVA 10: CIERRE Y ENTREGA ---
\begin{frame}[plain]
    \begin{tikzpicture}[remember picture,overlay]
        \fill[proesadark] (current page.south west) rectangle (current page.north east);
    \end{tikzpicture}
    
    \begin{center}
        \vspace{1.8cm}
        {\color{white}\Huge \textbf{Gracias}}
        
        \vspace{0.6cm}

        
        \vspace{1cm}
        {\color{white}\large \faEnvelope\ \ \ \href{mailto:contacto@proesa.org}{contacto@proesa.org}}
        
        \vspace{0.4cm}
        {\color{proesalight}\large \faGlobe\ \ \ \href{http://www.icesi.edu.co/proesa}{www.icesi.edu.co/proesa}}
    \end{center}
\end{frame}

\end{document}
